
\medskip
\phantomsection\label{cor:sga3-VIB-6.5.5}%
\noindent\textbf{Corollary 6.5.5.---}
\emph{Let \(G\) be an \(S\)-group scheme locally of finite type
and quasi-separated over \(S\), let \(H\) be a smooth group
scheme of finite presentation over \(S\) with connected fibers,
and let \(i:H\to G\) be a monomorphism of \(S\)-groups. Then:}
\begin{enumerate}
\item[\textup{a)}]\phantomsection\label{concl:sga3-VIB-6.5.5-a}
  \emph{\(\underline{\operatorname{Centr}}_G(H)\) and
  \(\underline{\operatorname{Norm}}_G(H)\) are representable by
  closed subschemes of \(G\), of finite presentation over \(G\).}
\item[\textup{a\(^{\prime}\))}]
  \phantomsection\label{concl:sga3-VIB-6.5.5-aprime}
  \emph{Likewise, for every monomorphism \(j:K\to G\) of
  \(S\)-group schemes that is of finite presentation,
  \(\underline{\operatorname{Transp}}_G(H,K)\) is representable
  by a closed subscheme of \(G\), of finite presentation over
  \(G\).}
\item[\textup{b)}]\phantomsection\label{concl:sga3-VIB-6.5.5-b}
  \emph{If \(G\) is quasi-compact, there exists an integer
  \(n\geq 0\) such that, if \(H_n\) denotes the \(n\)-th
  infinitesimal neighborhood of the unit section of \(H\), then}
  \[
    \underline{\operatorname{Centr}}_G(H)
      =\underline{\operatorname{Centr}}_G(H_n),
    \qquad
    \underline{\operatorname{Norm}}_G(H)
      =\underline{\operatorname{Norm}}_G(H_n),
  \]
  \[
    \underline{\operatorname{Transp}}_G(H,K)
      =\underline{\operatorname{Transp}}_G(H_n,K)
      =\underline{\operatorname{Transp}}_G(H_n,K_n).
  \]
\end{enumerate}

\emph{Proof.}%
\footnote{\label{note:sga3-VIB-66}Editor's note: the argument that
follows has been expanded. This also completes the insertion of
XI, \SGARefLink{sga3:xi:theorem:6:8}{6.8}--6.11; at \SGARefLink{def:sga3-VIB-6.6}{6.6} below we return to the original text of
VI\textsubscript{B}.}
First observe that the hypothesis on \(G\) implies that the
monomorphism \(H\to G\) is of finite presentation
(EGA~IV\(_1\), 1.2.4 and 1.4.3), as is the diagonal immersion
\[
  \Delta_{G/S}:G\longrightarrow G\times_S G
\]
(loc.\ cit., 1.4.3.1). The case of
\(\underline{\operatorname{Norm}}_G(H)\) therefore reduces to
that of the transporter, by taking \(H=K\).

In view of
\hyperref[ex:sga3-VIB-6.2.4-e]{6.2.4\textup{(e)}}, we apply
\hyperref[cor:sga3-VIB-6.5.4]{6.5.4} to the group scheme
\[
  X=H_G=H\times_S G
\]
over the base scheme \(G\), and to the following group subscheme
\(Y\).

In the case of
\(\underline{\operatorname{Transp}}_G(H,K)\), take \(Y\) to be
the inverse image of \(K_G\) under the morphism of \(G\)-groups
\[
  H\times_S G\longrightarrow G\times_S G,
  \qquad
  (h,g)\longmapsto(ghg^{-1},g).
\]
In the case of
\(\underline{\operatorname{Centr}}_G(H)\), take \(Y\) to be the
inverse image of the diagonal subgroup
\[
  \Delta_{G/S}(G)_G
\]
of \((G\times_S G)_G\) under the morphism of \(G\)-groups
\[
  H\times_S G\longrightarrow G\times_S G\times_S G,
  \qquad
  (h,g)\longmapsto(h,ghg^{-1},g).
\]
